% preamble-exam-zh.tex — 共享 preamble
% 用于所有 exam-zh 模板
%
% 样式策略：
%   屏幕 → 语义彩色（蓝=关键想法，红=易错，绿=路线，灰=提示/教师备注）
%   打印 → 自动降级为黑白（通过 print mode 或 monochrome 选项）

\documentclass{exam-zh}

% ---------- 基础配置 ----------
\examsetup{
  page/size = a4paper,
  page/show-head = false,
  page/show-foot = false,
  sealline/show = false,
  font = times,
  math-font = xits,
}

\usepackage{tcolorbox}
\usepackage{needspace}
\usepackage{graphicx}
\tcbuselibrary{skins,breakable}

% ---------- 打印/屏幕颜色切换 ----------
% 用法：\documentclass[monochrome]{exam-zh} 即可切换为纯黑白
% 注意：不要使用 \if@xxx 放在 makeatother 外部；这里改成普通条件名。
\newif\ifeduprintcolor
\eduprintcolortrue

\makeatletter
\@ifclasswith{exam-zh}{monochrome}{\eduprintcolorfalse}{}
\makeatother

% ---------- 语义颜色定义 ----------
% 屏幕：有色彩区分；打印：降级为灰度
\ifeduprintcolor
  % --- 屏幕彩色模式 ---
  \definecolor{edu-blue}{RGB}{47, 82, 122}       % 关键想法
  \definecolor{edu-blue-bg}{RGB}{243, 246, 252}
  \definecolor{edu-red}{RGB}{180, 40, 40}         % 易错提醒
  \definecolor{edu-red-bg}{RGB}{255, 245, 240}
  \definecolor{edu-green}{RGB}{34, 100, 60}       % 路线图
  \definecolor{edu-green-bg}{RGB}{242, 250, 245}
  \definecolor{edu-orange}{RGB}{160, 90, 0}       % 训练目标
  \definecolor{edu-orange-bg}{RGB}{255, 248, 235}
  \definecolor{edu-gray}{RGB}{100, 100, 100}      % 提示/教师备注
  \definecolor{edu-gray-bg}{RGB}{248, 248, 248}
  \definecolor{edu-purple}{RGB}{90, 50, 120}      % 思考题
  \definecolor{edu-purple-bg}{RGB}{248, 244, 252}
\else
  % --- 打印黑白模式 ---
  \definecolor{edu-blue}{gray}{0.15}
  \definecolor{edu-blue-bg}{gray}{0.96}
  \definecolor{edu-red}{gray}{0.25}
  \definecolor{edu-red-bg}{gray}{0.97}
  \definecolor{edu-green}{gray}{0.20}
  \definecolor{edu-green-bg}{gray}{0.97}
  \definecolor{edu-orange}{gray}{0.25}
  \definecolor{edu-orange-bg}{gray}{0.98}
  \definecolor{edu-gray}{gray}{0.40}
  \definecolor{edu-gray-bg}{gray}{0.97}
  \definecolor{edu-purple}{gray}{0.30}
  \definecolor{edu-purple-bg}{gray}{0.98}
\fi
\colorlet{edu-diagram-muted}{edu-gray}

% ---------- 自定义教学环境 ----------
% 共性：enhanced + breakable（跨页不断裂）

% 原题卡片 — 强边框，突出显示
\newtcolorbox{problemcard}{
  enhanced, breakable,
  colback=edu-blue-bg,
  colframe=edu-blue,
  boxrule=1.2pt,
  arc=0pt,
  left=3mm, right=3mm, top=3mm, bottom=3mm,
}

% 解题路线图 — 左边线 + 浅底
\newtcolorbox{routebox}{
  enhanced, breakable,
  colback=edu-green-bg,
  colframe=edu-green!30,
  boxrule=0pt,
  borderline west={2.5pt}{0pt}{edu-green},
  left=3mm, right=3mm, top=3mm, bottom=3mm,
}

% 关键想法 — 蓝色左边线
\newtcolorbox{keyideabox}{
  enhanced, breakable,
  colback=edu-blue-bg,
  colframe=edu-blue!30,
  boxrule=0pt,
  borderline west={2.5pt}{0pt}{edu-blue},
  left=3mm, right=3mm, top=3mm, bottom=3mm,
}

% 易错提醒 — 红色左边线
\newtcolorbox{mistakebox}{
  enhanced, breakable,
  colback=edu-red-bg,
  colframe=edu-red!20,
  boxrule=0pt,
  borderline west={3pt}{0pt}{edu-red},
  left=3mm, right=3mm, top=2.5mm, bottom=2.5mm,
}

% 提示 — 灰色虚线左边线
\newtcolorbox{hintbox}{
  enhanced, breakable,
  colback=edu-gray-bg,
  colframe=edu-gray!20,
  boxrule=0pt,
  borderline west={2pt}{0pt}{edu-gray, dashed},
  left=3mm, right=3mm, top=2mm, bottom=2mm,
}

% 思考题 — 紫色虚线左边线
\newtcolorbox{thinkbox}{
  enhanced, breakable,
  colback=edu-purple-bg,
  colframe=edu-purple!20,
  boxrule=0pt,
  borderline west={2pt}{0pt}{edu-purple, dashed},
  left=2.5mm, right=2.5mm, top=2.5mm, bottom=2.5mm,
}

% 教师备注 — 灰色左边线，小字
\newtcolorbox{teachernote}{
  enhanced, breakable,
  colback=edu-gray-bg,
  colframe=edu-gray!15,
  boxrule=0pt,
  borderline west={2pt}{0pt}{edu-gray!60},
  left=2.5mm, right=2.5mm, top=2.5mm, bottom=2.5mm,
  fontupper=\small,
  colupper=black!60,
}

% 训练目标 — 橙色左边线，小字
\newtcolorbox{traininggoal}{
  enhanced, breakable,
  colback=edu-orange-bg,
  colframe=edu-orange!15,
  boxrule=0pt,
  borderline west={1.5pt}{0pt}{edu-orange},
  left=2mm, right=2mm, top=1mm, bottom=1mm,
  fontupper=\small,
}

% 答题区 — 无边框，纯留白
\newcommand{\answerarea}[1][40mm]{%
  \vspace{2mm}%
  \begin{tcolorbox}[
    enhanced, breakable,
    colback=white,
    colframe=white,
    boxrule=0pt,
    height=#1,
    valign=top,
    bottom=0mm,
    after skip=0mm,
  ]
  \end{tcolorbox}%
}

\newcommand{\diagramcolfigure}[3]{%
  \begin{minipage}[t]{#2}
    \vspace{0pt}\centering
    \includegraphics[width=\linewidth]{\detokenize{#1}}
    \if\relax\detokenize{#3}\relax\else
      \par\vspace{1mm}{\scriptsize\color{edu-diagram-muted}#3}
    \fi
  \end{minipage}%
}

\newcommand{\diagramrowitem}[4]{%
  \begin{minipage}[t]{#2}
    \vspace{0pt}\centering
    \includegraphics[width=\linewidth]{\detokenize{#1}}
    \if\relax\detokenize{#3}\relax\else
      \par\vspace{1mm}{\scriptsize\bfseries #3}
    \fi
    \if\relax\detokenize{#4}\relax\else
      \par\vspace{0.5mm}{\scriptsize\color{edu-diagram-muted}#4}
    \fi
  \end{minipage}%
}

\newcommand{\answerareawithdiagramcol}[4]{%
  \vspace{2mm}%
  \begin{tcolorbox}[
    enhanced, breakable,
    colback=white,
    colframe=white,
    boxrule=0pt,
    height=#1,
    valign=top,
    bottom=0mm,
    after skip=0mm,
  ]
  \begin{minipage}[t]{\dimexpr\linewidth-#3-6mm\relax}
    \vspace{0pt}
  \end{minipage}\hfill
  \diagramcolfigure{#2}{#3}{#4}
  \end{tcolorbox}%
}

\newcommand{\answerareawithdiagram}[4]{%
  \answerareawithdiagramcol{#1}{#2}{#3}{#4}%
}

% ---------- 字体设置 ----------
% exam-zh (ctexbook) already sets CJK fonts via ctex.
% Only override if specific fonts are needed.
% macOS: Songti SC, STHeiti, STFangsong
% Linux: SimSun, SimHei, FangSong

% ---------- 页面设置 ----------
% exam-zh (ctexbook) already loads geometry.
% Override margins if needed:
% \geometry{top=16mm, bottom=16mm, left=14mm, right=14mm}

\examsetup{
  question/show-answer = true,
  fillin/show-answer = true,
  solution/show-solution = show-stay,
}

% ---------- 教师版解答样式 ----------
\newtcolorbox{solutionblock}{
  enhanced, breakable,
  colback=edu-blue-bg,
  colframe=edu-blue!25,
  boxrule=0pt,
  borderline west={2pt}{0pt}{edu-blue!50},
  arc=0pt,
  left=3mm, right=3mm, top=2mm, bottom=2mm,
  before skip=2mm,
  after skip=3mm,
  fontupper=\small,
}

% 解答步骤编号
\newcounter{solstep}
\newcommand{\solstep}[2]{%
  \stepcounter{solstep}%
  \par\medskip
  \noindent\hangindent=6mm
  {\color{edu-blue}\bfseries\textcircled{\scriptsize\thesolstep}\enspace #1}\enspace #2\par
}

\begin{document}

\section*{三角函数图像与三角形判定专题练习（教师版）}

\section*{一、选择题（每题 4 分，共 8 分）}


\needspace{14\baselineskip}
\begin{question}[points=4]
  \noindent
  \begin{minipage}[t]{0.64\linewidth}
    \vspace{0pt}
    函数 $f(x)=A\sin(\omega x+\varphi)$，$A>0,\omega>0,|\varphi|<\dfrac{\pi}{2}$。由图可知最大值为 $3$，$f(0)=\dfrac{3\sqrt2}{2}$，且峰值横坐标为 $\dfrac{\pi}{8}$，则 $f(x)$ 可以是（ ）
    \par
    \begin{choices}[columns=1]
      \item $3\sin\left(2x+\dfrac{\pi}{4}\right)$
      \item $3\sin\left(4x+\dfrac{\pi}{4}\right)$
      \item $\dfrac{3\sqrt2}{2}\sin\left(2x+\dfrac{\pi}{4}\right)$
      \item $3\sin\left(2x+\dfrac{\pi}{6}\right)$
    \end{choices}
  \end{minipage}\hfill
  \diagramcolfigure{diagram/jobs/c1-prompt/rendered/prompt.png}{0.30\linewidth}{读最大值、初值、峰值横坐标。}
  \paren[A]
  \begin{solutionblock}
    最大值给出 $A=3$；$3\sin\varphi=\dfrac{3\sqrt2}{2}$ 且 $|\varphi|<\dfrac{\pi}{2}$，所以 $\varphi=\dfrac{\pi}{4}$。峰值条件为 $\dfrac{\pi}{8}\omega+\dfrac{\pi}{4}=\dfrac{\pi}{2}$，得 $\omega=2$。
  \end{solutionblock}
\end{question}




\needspace{14\baselineskip}
\begin{question}[points=4]
  \noindent
  \begin{minipage}[t]{0.64\linewidth}
    \vspace{0pt}
    在 $\triangle ABC$ 中，$a,b,c$ 分别是角 $A,B,C$ 的对边。若 $B=\dfrac{\pi}{3}$，且 $a,b,c$ 成等差数列，则 $\triangle ABC$ 是（ ）
    \par
    \begin{choices}[columns=1]
      \item 等边三角形
      \item 只有 $AB=BC$ 的等腰三角形
      \item 直角三角形
      \item 不能确定形状
    \end{choices}
  \end{minipage}\hfill
  \diagramcolfigure{diagram/jobs/c2-prompt/rendered/prompt.png}{0.28\linewidth}{对照角 $B$ 与对边 $b$。}
  \paren[A]
  \begin{solutionblock}
    由等差数列得 $2b=a+c$。由余弦定理，$b^2=a^2+c^2-2ac\cos\dfrac{\pi}{3}=a^2+c^2-ac$。联立可得 $3(a-c)^2=0$，所以 $a=c$，进而 $a=b=c$。
  \end{solutionblock}
\end{question}



\section*{二、填空题（每题 4 分，共 8 分）}


\begin{question}[points=4]
  函数 $f(x)=A\sin(\omega x+\varphi)$ 的最大值为 $4$，$f(0)=2$，且在 $x=\dfrac{\pi}{12}$ 处达到峰值。若 $0<\omega<6$，则 $\omega=$ \underline{\hspace{1.8cm}}。
  \paren[$4$]
  \begin{solutionblock}
    $A=4$，$4\sin\varphi=2$，所以 $\varphi=\dfrac{\pi}{6}$。峰值条件：$\dfrac{\pi}{12}\omega+\dfrac{\pi}{6}=\dfrac{\pi}{2}$，得 $\omega=4$。
  \end{solutionblock}
\end{question}



\begin{question}[points=4]
  在 $\triangle ABC$ 中，若 $B=\dfrac{\pi}{3}$，且 $a,b,c$ 成等差数列，则 $\triangle ABC$ 的形状是 \underline{\hspace{2.4cm}}。
  \paren[等边三角形]
  \begin{solutionblock}
    同选择题第 2 题，$2b=a+c$ 与 $b^2=a^2+c^2-ac$ 联立，推出 $a=b=c$。
  \end{solutionblock}
\end{question}



\needspace{15\baselineskip}
\begin{center}
  \diagramrowitem{diagram/jobs/f1-prompt/rendered/prompt.png}{0.24\linewidth}{第 1 题}{}\hfill  \diagramrowitem{diagram/jobs/f2-prompt/rendered/prompt.png}{0.22\linewidth}{第 2 题}{}\end{center}


\section*{三、解答题（14 分）}

\needspace{8\baselineskip}

\begin{problem}[points=14]
  已知函数 $f(x)=A\sin(\omega x+\varphi)$，其中 $A>0,\omega>0,|\varphi|<\dfrac{\pi}{2}$。图像显示最大值为 $3$，$f(0)=\dfrac32$，并且在 $x=\dfrac{\pi}{6}$ 处达到峰值，且 $0<\omega<3$。
\begin{enumerate}[label=(\arabic*)]
  \item 求 $f(x)$ 的解析式；
  \item 在 $\triangle ABC$ 中，$a,b,c$ 分别是角 $A,B,C$ 的对边。若 $f(B)=\dfrac32$，且 $a,b,c$ 成等差数列，判断 $\triangle ABC$ 的形状。
\end{enumerate}

  \noindent\textbf{(1)}\par
  \answerareawithdiagramcol{35mm}{diagram/jobs/p1-part1-prompt/rendered/prompt.png}{0.32\linewidth}{读图求参。}
  \noindent\textbf{(2)}\par
  \answerareawithdiagramcol{40mm}{diagram/jobs/p1-part2-prompt/rendered/prompt.png}{0.32\linewidth}{对照角 $B$ 与对边 $b$。}
  \setcounter{solstep}{0}
  \begin{solutionblock}
    \textbf{答案：}(1) $f(x)=3\sin\left(2x+\dfrac{\pi}{6}\right)$；(2) $\triangle ABC$ 为等边三角形。\par\medskip
    \solstep{读图求 $A$ 与 $\varphi$}{最大值给出 $A=3$；由 $3\sin\varphi=\dfrac32$ 且 $|\varphi|<\dfrac{\pi}{2}$，得 $\varphi=\dfrac{\pi}{6}$。}
    \solstep{用峰值求 $\omega$}{$\dfrac{\pi}{6}\omega+\dfrac{\pi}{6}=\dfrac{\pi}{2}+2k\pi$，所以 $\omega=2+12k$。由 $0<\omega<3$，得 $\omega=2$。}
    \solstep{由函数值求角 $B$}{$3\sin\left(2B+\dfrac{\pi}{6}\right)=\dfrac32$，所以 $B=\dfrac{\pi}{3}$。}
    \solstep{判定三角形}{$2b=a+c$，且 $b^2=a^2+c^2-ac$。联立得 $3(a-c)^2=0$，所以 $a=c$，进一步 $a=b=c$。}
  \end{solutionblock}
\end{problem}


\clearpage
\section*{参考答案}


\begin{keyideabox}
  \textbf{选择题}
  
\end{keyideabox}



\begin{keyideabox}
  \textbf{填空题}
  
\end{keyideabox}



\begin{keyideabox}
  \textbf{解答题}
  
\end{keyideabox}



\end{document}
